% Draft grounded in repository exports; methodology.tex is authoritative.
% Audit trail: evaluation_sources.md in this directory. No figures generated.
\chapter{Evaluation \& Results}
\label{ch:evaluation}
This chapter evaluates the three experimental tracks defined in Chapter~\ref{ch:methodology}: supervised \gls{OD}, zero-shot prompt-based localisation, and multi-attribute traffic-sign classification. The analysis considers predictive performance alongside model capacity, sensitivity to the experimental configuration and computational requirements. Results are interpreted within the two Maltese street-level domains, with particular attention to small-object localisation, semantic ambiguity and the recognition of physical condition.

The chapter draws on the retained test summaries, prediction exports and evaluation records. Where these records differ from the protocol established in Chapter~\ref{ch:methodology}, the discrepancy is identified rather than treating the available measurement as an exact implementation of that protocol. Validation-only evaluations, earlier attribute experiments and prompt-localisation pilots are not substituted for final test results. All predictive measures in the tables are expressed as percentages unless otherwise stated; differences between them are expressed in percentage points.

\section{Evaluation Framework and Performance Measures}
\label{sec:ev-framework}
\subsection{Evaluation Protocol and Ground-Truth Matching}
\label{sec:ev-protocol}
The held-out localisation sets comprise 369 \gls{MDWD} images and 747 \gls{MTSD} images. Attribute classification instead uses the 1,890 eligible test crops inherited from the source-image partitions, following the minimum-size restriction in Section~\ref{sec:attr-cls}. The classification results therefore quantify characterisation of annotated objects with known bounding boxes, rather than the performance of a complete detection-and-characterisation pipeline.

For localisation, a prediction constitutes a \gls{TP} when it is assigned to an eligible \gls{GT} box at \(\mathrm{IoU}\geq0.50\), subject to one-to-one assignment. Unmatched predictions constitute \glspl{FP}, while unmatched eligible annotations constitute \glspl{FN}. Supervised evaluation additionally requires agreement with the annotated category. In prompt-based evaluation, eligibility is determined by the query-specific target-class mapping in Section~\ref{sec:prompt-matching}. Images containing no eligible targets remain in the evaluation, so a model cannot obtain high precision merely by being assessed on scenes known to contain the requested object.

The retained prompt evaluator processes predictions in descending score order and assigns each prediction to its highest-overlap target. A further prediction whose highest-overlap target has already been assigned is counted as a duplicate and a \gls{FP}; the implementation does not seek a second-best unassigned target in that case. This detail is relevant in crowded scenes and distinguishes the stored implementation from an optimal global assignment. Predictions from LocateAnything and Cosmos Reason2 have uniform scores, preserving generation order. The confidence threshold of 0.30, minimum box area of 100 pixels$^2$, \gls{NMS} threshold of 0.50 and dataset-specific output limits follow Table~\ref{tab:prompt-fixed-settings}.

Training and validation remain distinct from test evaluation. In particular, the attribute checkpoint is selected by validation mean macro-\(F_1\), and only the selected checkpoint is assessed on test crops. For \gls{MDWD}, the native summaries contain separate validation and test fields; only the explicitly identified test fields are used in the principal detector table. Their precision and recall retain the framework's native reporting convention, and the stored summary \(F_1\) is retained without treating it as a pooled object-level count statistic.

A further distinction is necessary for the unified \gls{MTSD} exports. Their precision, recall and standard \(F_1\) are calculated from pooled \gls{TP}, \gls{FP} and \gls{FN} counts using Equation~\ref{eq:ev-prf}, at the confidence value maximising \(F_1\) on the evaluated partition. The test tables report this ordinary harmonic-mean \(F_1\) without a separate starred notation. Because the operating threshold is selected on the same test partition, however, the reported value describes the best stored test-set precision--recall trade-off and is not an unbiased deployment estimate. Confidence-ranked \gls{AP} remains the principal measure for supervised configuration comparisons. The unified YOLO export retains predictions down to confidence 0.001, with \gls{NMS} at 0.70 and a maximum of 100 detections, so its prediction-generation settings also differ from the prompt pipeline.

\paragraph{MTSD taxonomy and implementation limitation.}
Methodology excludes \emph{Tourist Sign} from the experiment-ready localisation taxonomy and specifies that its boxes remain as ignore regions. The available unified exports evaluate 11 categories, containing 1,975 target boxes across the 747 test images. However, the retained supervised evaluator removes the excluded category's annotations and predictions by category identifier, without suppressing otherwise eligible predictions that overlap its boxes. The retained prompt matching code likewise contains no explicit Tourist Sign overlap-ignore operation. Moreover, some supervised run records explicitly retain the training taxonomy while excluding Tourist Sign only during evaluation. The stored measurements therefore cannot be described as fully implementing the specified exclusion-and-ignore protocol. The \gls{MTSD} localisation tables below report these available exports with this qualification; no numerical correction is inferred from the three excluded test instances. Their effect on precision or ranking remains unquantified.

\subsection{Object Detection and Localisation Metrics}
\label{sec:ev-detection-metrics}
Bounding-box overlap is measured by \gls{IoU}, as defined in Equation~\ref{eq:kappa-iou}. At the stated matching and confidence thresholds, precision, recall and \(F_1\) are defined as:
\begin{equation}
\label{eq:ev-prf}
P=\frac{\mathrm{TP}}{\mathrm{TP}+\mathrm{FP}},\qquad
R=\frac{\mathrm{TP}}{\mathrm{TP}+\mathrm{FN}},\qquad
F_1=\frac{2PR}{P+R}.
\end{equation}
\noindent Precision describes the reliability of reported objects, while recall measures annotation coverage. Their harmonic mean penalises configurations for which one is high and the other low. These quantities are calculated from pooled counts in the unified detector evaluator, whereas prompt results are first calculated independently for each query and then macro-averaged over the eligible prompt inventory. The macro-average of prompt \(F_1\) is not, in general, the harmonic mean of macro-averaged precision and recall.

\gls{AP}@50 summarises the confidence-ranked precision--recall curve at \(\mathrm{IoU}=0.50\). Averaging category-level values gives \gls{mAP}@50, while \gls{mAP}@50--95 additionally averages across the ten overlap thresholds from 0.50 to 0.95 in increments of 0.05. The latter measure places greater emphasis on accurate box extent. The unified evaluator also retains \gls{AP}@75 and size-stratified \gls{AP} and average recall (AR). The small, medium and large area ranges use the conventional boundaries of \(32^2\) and \(96^2\) pixels$^2$ in the annotation coordinate system. They must not be confused with the relative-area statistics of the source corpora or with size bins calculated after resizing to 640 pixels.

For cross-model prompt localisation, precision, recall, \(F_1\), mean matched-box \gls{IoU} and runtime are the common measures. Confidence-ranked \gls{AP} is not used as a universal headline metric because LocateAnything and Cosmos Reason2 do not provide meaningful per-box confidence scores. Mean matched-box \gls{IoU} is conditional on successful matching: a model may tightly localise a small subset of objects while missing most eligible instances. It is therefore interpreted jointly with recall rather than as a measure of overall localisation quality.

\subsection{Attribute-Classification Metrics}
\label{sec:ev-attribute-metrics}
Each attribute head is evaluated using accuracy, macro-\(F_1\), per-class \(F_1\) and a confusion matrix. Accuracy is the fraction of test crops assigned the correct attribute label. For attribute \(a\), macro-\(F_1\) gives equal weight to its \(C_a\) classes:
\begin{equation}
F_{1,a}^{\mathrm{macro}}=\frac{1}{C_a}\sum_{c=1}^{C_a}F_{1,a,c},\qquad
\overline{F_1}^{\mathrm{macro}}=\frac{1}{4}\sum_{a=1}^{4}F_{1,a}^{\mathrm{macro}}.
\end{equation}
\noindent The second quantity is the principal representation-comparison measure and gives equal weight to viewing angle, mounting type, physical condition and geometric shape. This avoids allowing the most frequent labels or the attribute with the largest class inventory to dominate the overall assessment. Macro-averaging is particularly relevant for physical condition, where \emph{Good} signs substantially outnumber \emph{Heavily Damaged} signs. Confusion matrices retain the direction of errors, distinguishing, for example, a damaged sign assigned a less severe state from a good sign incorrectly flagged for deterioration.

\subsection{Computational-Efficiency Measures}
\label{sec:ev-efficiency-measures}
Model capacity is described through total parameters and, where an explicit audit exists, trainable parameters. Serialized checkpoint size is reported separately because it depends on numerical precision and saved content as well as architecture. The MDWD parameter counts retain the run-specific audits and need not equal the rounded pretrained-checkpoint counts in Table~\ref{tab:model-matrix}. A corresponding complete run-specific parameter and file-size audit was not located for the MTSD detector exports, so the pretrained capacity table supplies architectural context rather than substituted measurements. A small trainable fraction reduces the parameters updated during adaptation but does not remove the frozen backbone from inference.

Latency and throughput retain the boundaries of the corresponding measurement. Native \gls{MDWD} test timings, the \gls{MTSD} unified prediction loop, prompt prediction-pipeline latency and batched attribute evaluation do not time identical operations. Prompt latency follows Section~\ref{sec:inference-execution} and applies to one image--prompt pair. The \gls{MTSD} unified detector loop includes image loading, prediction and conversion of detections into export records. Attribute timing distinguishes the model forward pass from end-to-end crop evaluation, with three timing repetitions after five warm-up batches. Where FPS is obtained from latency, \(\mathrm{FPS}=1000/t_{\mathrm{ms}}\); such a reciprocal is not a measured multi-query or complete municipal-system throughput. Training durations are reported where retained, while missing model-size or memory audits are not replaced with theoretical estimates.

\subsection{Statistical Uncertainty}
\label{sec:ev-uncertainty}
The retained uncertainty analyses use percentile bootstrap intervals with 2,000 resamples and seed 42. Table~\ref{tab:ev-mdwd-ci} reports an image-level bootstrap for the stored YOLO26-L \gls{MDWD} test predictions at confidence 0.25. Images, rather than individual boxes, form the resampling unit, preserving the dependence among objects within a scene. These intervals apply to the audited operating point and must not be attached to the different native summary scores.

% SOURCE: Documents/Final-Reports/Statistical-Uncertainty/20260908-105836/bootstrap_results.csv and bootstrap_config.json; yolo26l@YOLO26-EUVIP test only.
\begin{table}[htbp]
\centering
\small
\caption{Stored image-bootstrap intervals for YOLO26-L on the 369-image \gls{MDWD} test split at confidence 0.25. Values are percentages; 2,000 resamples, seed 42.}
\label{tab:ev-mdwd-ci}
\renewcommand{\arraystretch}{1.15}
\begin{tabularx}{\textwidth}{@{}>{\raggedright\arraybackslash}Xrr@{}}
\toprule
Measure & Estimate & 95\% interval \\
\midrule
Precision & 94.54 & [92.53, 96.28] \\
Recall & 87.61 & [84.17, 90.73] \\
$F_1$ & 90.94 & [88.46, 93.01] \\
Matched IoU & 92.15 & [91.51, 92.78] \\
\bottomrule
\end{tabularx}
\end{table}


The refreshed uncertainty export also provides 95\% intervals for the retained 1,890-crop attribute snapshot and image-level intervals for the final full-split PromptDetect runs on both datasets. Attribute resampling reconstructs test-crop outcomes from the stored confusion matrices and therefore retains an independent-crop approximation; it does not cluster different crops from a common source image. PromptDetect resampling uses the image as the unit and preserves all prompt-specific outcomes for that image. These analyses quantify sampling uncertainty in the retained predictions, but they do not capture variation across training seeds or establish geographic generalisation.

No paired comparison of the leading current attribute configurations was located. Small differences between those models are therefore reported descriptively, without claims of statistical significance; overlapping marginal intervals alone are not a paired test. The PromptDetect intervals similarly quantify each retained model--prompt result and do not remove the systematic effects of prompt selection, output post-processing or model-specific inference settings.

\clearpage
\section{Supervised Object Detection Results}
\label{sec:ev-supervised}
\subsection{MDWD Cross-Architecture and Model-Scale Results}
\label{sec:ev-mdwd-supervised}
Table~\ref{tab:ev-mdwd-benchmark} compares the retained YOLO and RF-DETR configurations on the held-out \gls{MDWD} test set. The RF-DETR-N run is supported by the retained local training artefacts, while the RF-DETR-S and RF-DETR-M measurements were exported from the completed Roboflow training runs and supplied in the final cross-architecture results table. All rows report the standard precision, recall and harmonic-mean \(F_1\) defined in Equation~\ref{eq:ev-prf}.

% SOURCE: Results/MDWD-Results/{YOLO11,YOLO12,YOLO26}/Model-Size-Comparison/*_summary.csv, explicit test_ columns; final MDWD cross-architecture Table IV supplied by the researcher for RF-DETR-N/S/M Roboflow results.
\begin{table}[htbp]
\centering
\small
\caption{\gls{MDWD} held-out test results from the final cross-architecture benchmark. Parameters are in millions and checkpoint size in MB; predictive measures are percentages.}
\label{tab:ev-mdwd-benchmark}
\renewcommand{\arraystretch}{1.15}
\begin{tabularx}{\textwidth}{@{}>{\raggedright\arraybackslash}Xrrrrrrr@{}}
\toprule
Model & Params. & MB & \gls{mAP}@50 & \gls{mAP}@50--95 & $P$ & $R$ & $F_1$ \\
\midrule
YOLO11-N & 2.59 & 5.22 & 85.71 & 67.13 & 92.56 & 78.08 & 84.71 \\
YOLO11-S & 9.43 & 18.29 & 90.06 & 73.34 & 94.24 & 85.14 & 89.46 \\
YOLO11-M & 20.06 & 38.64 & 89.72 & 74.77 & 94.19 & 85.68 & 89.73 \\
YOLO12-N & 2.57 & 5.26 & 88.01 & 70.29 & 94.02 & 80.39 & 86.67 \\
YOLO12-S & 9.26 & 18.06 & 90.00 & 74.41 & 95.72 & 85.13 & 90.11 \\
YOLO12-M & 20.14 & 38.88 & 90.30 & 75.84 & 95.72 & 85.73 & 90.45 \\
YOLO26-N & 2.51 & 5.14 & 86.59 & 68.86 & 93.62 & 79.06 & 85.73 \\
YOLO26-S & 9.95 & 19.38 & 90.14 & 75.27 & 94.89 & 86.19 & 90.33 \\
YOLO26-M & 21.78 & 42.00 & 89.51 & 75.97 & 93.07 & 86.25 & 89.53 \\
YOLO26-L & 26.18 & 50.54 & 89.76 & 78.20 & 97.18 & 86.98 & 91.80 \\
RF-DETR-N & 30.16 & 115.24 & 90.21 & 70.88 & 95.17 & 85.53 & 90.10 \\
RF-DETR-S & 31.81 & 121.52 & 92.61 & 76.17 & 96.39 & 89.24 & 92.68 \\
RF-DETR-M & 33.38 & 127.54 & 94.49 & 78.21 & 96.63 & 90.69 & 93.56 \\
\bottomrule
\end{tabularx}
\end{table}


RF-DETR-M achieves the strongest overall MDWD result, with test \gls{mAP}@50 of 94.49\%, recall of 90.69\% and \(F_1\) of 93.56\%. Its \gls{mAP}@50--95 of 78.21\% is effectively tied with YOLO26-L at 78.20\%, while YOLO26-L retains the highest precision at 97.18\%. RF-DETR-S is also the strongest medium-sized configuration by \(F_1\), reaching 92.68\% compared with 90.45\% for YOLO12-M, 89.73\% for YOLO11-M and 89.53\% for YOLO26-M. The cross-architecture result therefore favours RF-DETR for balanced retrieval, while the largest YOLO configuration remains competitive for strict box localisation and offers a substantially smaller checkpoint.

Increasing model scale generally improves \gls{mAP}@50--95, although the gains are not proportional to parameter count. YOLO26-S reaches 75.27\% with 9.95 million parameters, whereas YOLO26-M reaches 75.97\% with 21.78 million. The 0.70-point improvement therefore accompanies more than twice as many parameters. YOLO12-S and YOLO12-M improve from 74.41\% to 75.84\%, while YOLO11-S and YOLO11-M improve from 73.34\% to 74.77\%. The larger configurations may be justified where tighter localisation is valuable, but the small models already recover much of the available performance.

The RF-DETR scale trend is clearer in the standard \(F_1\) measure: RF-DETR-N, -S and -M increase from 90.10\% to 92.68\% and 93.56\%, respectively. This improvement comes with comparatively large serialized checkpoints, ranging from 115.24 to 127.54\,MB. RF-DETR-M has 33.38 million parameters and a 127.54\,MB checkpoint, compared with 26.18 million and 50.54\,MB for YOLO26-L, so its higher recall and \(F_1\) require a larger storage footprint.

The class-level audit in Table~\ref{tab:ev-mdwd-classes} clarifies what the aggregate scores conceal. At confidence 0.25, YOLO26-L achieves its highest class \(F_1\) on \emph{Organic Waste}, followed by \emph{Recyclable Material} and \emph{Mixed Waste}. \emph{Other Waste} has lower precision and recall, consistent with the heterogeneous visual scope of this residual category. \emph{Orange CMD} retains high precision but lower recall despite its distinctive appearance. Its comparatively small support and object scale provide plausible contributing factors, although the aggregate class metrics do not isolate their individual effects.

% SOURCE: Documents/Final-Reports/Robustness-Slices/20260908-103240/robustness_slice_results.csv; YOLO26-EUVIP, test, class slices.
\begin{table}[htbp]
\centering
\small
\caption{Class-level \gls{MDWD} results for YOLO26-L, recomputed in the stored audit at confidence 0.25 and matching \gls{IoU} 0.50. Measures are percentages; this operating point differs from Table~\ref{tab:ev-mdwd-benchmark}.}
\label{tab:ev-mdwd-classes}
\renewcommand{\arraystretch}{1.15}
\begin{tabularx}{\textwidth}{@{}>{\raggedright\arraybackslash}Xrrrrr@{}}
\toprule
Class & Support & $P$ & $R$ & $F_1$ & Matched \gls{IoU} \\
\midrule
Mixed Waste & 334 & 95.82 & 89.22 & 92.40 & 91.83 \\
Orange CMD & 78 & 93.85 & 78.21 & 85.31 & 94.09 \\
Organic Waste & 178 & 97.13 & 94.94 & 96.02 & 93.31 \\
Other Waste & 278 & 91.63 & 78.78 & 84.72 & 89.88 \\
Recyclable Material & 238 & 94.07 & 93.28 & 93.67 & 93.40 \\
\bottomrule
\end{tabularx}
\end{table}


% FIGURE PLACEHOLDER: Plot MDWD test F1 and mAP@50--95 against parameter count for the ten retained YOLO variants and three RF-DETR variants. Distinguish the locally retained YOLO summaries from the researcher-supplied Roboflow RF-DETR exports.

\subsection{MTSD Cross-Architecture and Model-Scale Results}
\label{sec:ev-mtsd-supervised}
The complete 13-configuration unified export provides the broadest available \gls{MTSD} scale comparison under the mild training-data strategy. Table~\ref{tab:ev-mtsd-benchmark} preserves this common block, rather than mixing the highest-scoring configuration of each family across different augmentation strategies and resolutions. RF-DETR uses its scale-specific native input resolutions, so its scale comparison includes changes in spatial input as well as capacity. The evaluation and training-taxonomy qualifications in Section~\ref{sec:ev-protocol} apply throughout this subsection.

% SOURCE: Results/MTSD-Results/Unified-Evaluation-NoTourist/20260722-all13/*/{metrics_summary,evaluation_record}.json.
\begin{table}[htbp]
\centering
\small
\caption{\gls{MTSD} mild-augmentation scale comparison on 747 test images and 11 evaluated categories, from the 22 July unified export. All measures except input size are percentages. $P$, $R$ and standard $F_1$ are reported at each run's best test-set confidence; the Tourist Sign ignore-region discrepancy described in Section~\ref{sec:ev-protocol} applies.}
\label{tab:ev-mtsd-benchmark}
\renewcommand{\arraystretch}{1.15}
\begin{tabularx}{\textwidth}{@{}>{\raggedright\arraybackslash}Xrrrrrr@{}}
\toprule
Model & Input & \gls{mAP}@50 & \gls{mAP}@50--95 & $P$ & $R$ & $F_1$ \\
\midrule
YOLO11-N & 640 & 60.16 & 54.38 & 88.34 & 54.84 & 67.67 \\
YOLO11-S & 640 & 63.44 & 57.54 & 89.95 & 57.97 & 70.50 \\
YOLO11-M & 640 & 65.86 & 60.14 & 91.61 & 59.75 & 72.33 \\
YOLO12-N & 640 & 60.07 & 54.66 & 90.28 & 54.53 & 67.99 \\
YOLO12-S & 640 & 63.79 & 57.80 & 88.74 & 58.28 & 70.35 \\
YOLO12-M & 640 & 63.55 & 57.79 & 85.51 & 58.89 & 69.75 \\
YOLO26-N & 640 & 60.19 & 54.38 & 87.48 & 53.42 & 66.33 \\
YOLO26-S & 640 & 65.88 & 58.98 & 86.70 & 59.39 & 70.49 \\
YOLO26-M & 640 & 67.17 & 61.26 & 88.55 & 61.11 & 72.32 \\
YOLO26-L & 640 & 68.63 & 62.01 & 85.34 & 63.09 & 72.55 \\
RF-DETR-N & 384 & 71.46 & 62.29 & 90.61 & 65.47 & 76.01 \\
RF-DETR-S & 512 & 77.28 & 66.84 & 93.02 & 68.20 & 78.70 \\
RF-DETR-M & 576 & 78.42 & 68.59 & 92.55 & 70.43 & 79.99 \\
\bottomrule
\end{tabularx}
\end{table}


RF-DETR-M leads this baseline block with \gls{mAP}@50--95 of 68.59\%, followed by RF-DETR-S at 66.84\% and RF-DETR-N at 62.29\%. The highest YOLO baseline is YOLO26-L at 62.01\%. The 6.58-point difference between RF-DETR-M and YOLO26-L is substantial descriptively, but it compares the complete trained configurations, including their pretrained initialisation, model-specific preprocessing and training procedures. It does not isolate architecture as the sole cause.

Within the YOLO families, the nano configurations are closely grouped near 54--55\% \gls{mAP}@50--95. Scaling improves YOLO11 and YOLO26, whereas YOLO12-S and -M are effectively tied at approximately 57.8\%. Additional capacity alone therefore does not consistently resolve traffic-sign detection difficulty at 640 pixels. The standard \(F_1\) results make the precision--recall imbalance explicit: even at each run's best test confidence, precision remains much higher than recall, showing that a substantial fraction of annotated signs is not recovered.

% SOURCE: Results/MTSD-Results/Unified-Evaluation-NoTourist/20260818-084730-pending-yolo-test/*/metrics_summary.json; Strong-640 entries only.
\begin{table}[htbp]
\centering
\small
\caption{Available strong-augmentation test exports at 640 pixels. Percentages; $F_1$ is the standard harmonic mean at the best test-set confidence. The same evaluation limitation as Table~\ref{tab:ev-mtsd-benchmark} applies.}
\label{tab:ev-mtsd-strong}
\renewcommand{\arraystretch}{1.15}
\begin{tabularx}{\textwidth}{@{}>{\raggedright\arraybackslash}Xrrrrr@{}}
\toprule
Model & \gls{mAP}@50 & \gls{mAP}@50--95 & $P$ & $R$ & $F_1$ \\
\midrule
YOLO11-N & 66.12 & 59.24 & 88.91 & 60.10 & 71.72 \\
YOLO11-S & 70.42 & 62.83 & 89.66 & 63.65 & 74.44 \\
YOLO11-M & 72.63 & 65.36 & 89.23 & 67.95 & 77.15 \\
YOLO26-N & 68.04 & 59.66 & 84.02 & 61.77 & 71.20 \\
YOLO26-S & 72.17 & 64.21 & 90.09 & 63.49 & 74.49 \\
YOLO26-M & 75.65 & 67.99 & 86.53 & 69.62 & 77.16 \\
\bottomrule
\end{tabularx}
\end{table}


The newer strong-augmentation exports alter the picture. At 640 pixels, YOLO26-M reaches 67.99\% \gls{mAP}@50--95, already approaching the mild RF-DETR-M baseline, while YOLO11-M reaches 65.36\%. These gains show why the baseline architectural ranking should not be treated as invariant to data preparation. The available strong test block is incomplete: several YOLO12 configurations and native-resolution RF-DETR-S/M strong configurations have validation exports without equivalent retained unified test outputs. A completed training run or a validation-only result is not counted as a final test measurement.

% SOURCE: Results/MTSD-Results/Unified-Evaluation-NoTourist/20260722-all13/rfdetr-m/metrics_summary.json; Results/MTSD-Runs/YOLO26-MTSD/strongaug-wandb-followup-yolo26m-img1280-s42/unified_evaluation/metrics_summary.json; per_class fields.
\begin{table}[htbp]
\centering
\small
\caption{Per-class \gls{MTSD} test \gls{AP}@50--95 (\%) for the strongest mild baseline and the highest-scoring available strong configuration. Different training strategies and input sizes preclude attributing the difference solely to architecture.}
\label{tab:ev-mtsd-classes}
\renewcommand{\arraystretch}{1.15}
\begin{tabularx}{\textwidth}{@{}>{\raggedright\arraybackslash}Xrrr@{}}
\toprule
Class & Support & RF-DETR-M & YOLO26-M \\
\midrule
Pedestrian Crossing & 143 & 71.65 & 85.04 \\
Stop Sign & 130 & 86.71 & 89.62 \\
No Entry & 211 & 82.78 & 87.93 \\
Roundabout Ahead & 67 & 76.96 & 85.16 \\
No Through Road & 61 & 92.58 & 92.77 \\
Blind-Spot Mirror & 160 & 83.62 & 88.05 \\
Street Sign & 94 & 63.59 & 68.44 \\
Directional Sign & 108 & 48.23 & 57.85 \\
Auxiliary Sign & 120 & 34.73 & 48.83 \\
Back-Unknown & 661 & 62.69 & 70.96 \\
Other-Unknown & 220 & 50.92 & 67.59 \\
\bottomrule
\end{tabularx}
\end{table}


The class-level comparison in Table~\ref{tab:ev-mtsd-classes} contrasts the mild RF-DETR-M baseline with the highest-scoring available strong configuration, YOLO26-M at 1280 pixels. \emph{No Through Road} exceeds 92\% \gls{AP}@50--95 in both cases, while \emph{Auxiliary Sign}, \emph{Directional Sign} and \emph{Other-Unknown} remain much weaker. Auxiliary Sign improves from 34.73\% to 48.83\%, yet remains the least accurately detected category in the latter configuration. Class frequency alone does not explain this ordering: the 61 No Through Road instances are detected more accurately than the 120 Auxiliary Sign instances. Visual scale, semantic variability and the available discriminative detail are therefore important alongside support.

% SOURCE: Results/MTSD-Results/Unified-Evaluation-NoTourist/20260722-all13/rfdetr-m/metrics_summary.json; 20260818-084730-pending-yolo-test/YOLO26M-Strong-{640,960}/metrics_summary.json; Results/MTSD-Runs/YOLO26-MTSD/strongaug-wandb-followup-yolo26m-img1280-s42/unified_evaluation/metrics_summary.json.
\begin{table}[htbp]
\centering
\small
\caption{\gls{MTSD} size-stratified \gls{AP} and AR (\%) over \gls{IoU} 0.50--0.95. AR permits up to 100 detections; size bins use the annotation coordinate system, not resized detector pixels.}
\label{tab:ev-size-ap}
\renewcommand{\arraystretch}{1.15}
\begin{tabularx}{\textwidth}{@{}>{\raggedright\arraybackslash}Xrrrrrr@{}}
\toprule
Configuration & $AP_S$ & $AP_M$ & $AP_L$ & $AR_S$ & $AR_M$ & $AR_L$ \\
\midrule
RF-DETR-M, mild, 576 & 13.76 & 39.10 & 87.36 & 22.09 & 51.02 & 92.45 \\
YOLO26-M, strong, 640 & 12.19 & 44.25 & 83.07 & 19.38 & 56.15 & 86.84 \\
YOLO26-M, strong, 960 & 21.03 & 55.64 & 86.21 & 33.39 & 66.07 & 89.07 \\
YOLO26-M, strong, 1280 & 33.41 & 64.04 & 86.52 & 49.95 & 73.35 & 90.98 \\
\bottomrule
\end{tabularx}
\end{table}


Size-stratified results in Table~\ref{tab:ev-size-ap} provide direct evidence of the small-sign limitation. For strong YOLO26-M, small-object \gls{AP} increases from 12.19\% at 640 pixels to 33.41\% at 1280 pixels, whereas large-object \gls{AP} increases from 83.07\% to 86.52\%. The much larger gain for small objects supports the interpretation that additional spatial detail is particularly valuable for distant signs. Nevertheless, small-object AR remains below 50\% even at 1280 pixels, indicating that the highest aggregate score does not imply reliable coverage of the smallest eligible signs.

\subsection{Effect of MTSD Training-Data Augmentation}
\label{sec:ev-augmentation}
The none, mild and strong training-data strategies are compared in Table~\ref{tab:ev-augmentation}, retaining a common 640-pixel input wherever test results are available. The strong strategy improves all available comparisons against mild augmentation. YOLO26-M increases from 61.26\% to 67.99\% \gls{mAP}@50--95, while YOLO11-M increases from 60.14\% to 65.36\%. For the small variants, YOLO11-S and YOLO26-S gain 5.29 and 5.22 points respectively over their mild baselines.

% SOURCE: Results/MTSD-Results/Unified-Evaluation-NoTourist/{20260722-all13,20260818-084730-pending-yolo-test}/*/metrics_summary.json; Results/MTSD-Runs/*/*noaug_img640*/unified_evaluation/metrics_summary.json.
\begin{table}[htbp]
\centering
\small
\caption{Available \gls{MTSD} test \gls{mAP}@50--95 (\%) at 640 pixels by training-data strategy. These are configuration comparisons: some earlier runs differ in effective batch size, so the table does not establish an isolated causal augmentation effect.}
\label{tab:ev-augmentation}
\renewcommand{\arraystretch}{1.15}
\begin{tabularx}{\textwidth}{@{}>{\raggedright\arraybackslash}Xrrr@{}}
\toprule
Model & None & Mild & Strong \\
\midrule
YOLO11-M & 59.16 & 60.14 & 65.36 \\
YOLO12-M & 57.25 & 57.79 & -- \\
YOLO26-M & 61.57 & 61.26 & 67.99 \\
YOLO11-S & 57.01 & 57.54 & 62.83 \\
YOLO26-S & 58.34 & 58.98 & 64.21 \\
\bottomrule
\end{tabularx}
\end{table}


Mild augmentation is not uniformly better than no offline augmentation. YOLO26-M records 61.57\% without augmentation and 61.26\% with mild augmentation, whereas YOLO11-M improves from 59.16\% to 60.14\%. This pattern argues against treating augmentation intensity as a monotonic performance control. The strong strategy produces the clearer improvements in the available records, but the comparison should be interpreted at configuration level: the historical and later YOLO runs do not all share the same effective optimiser batch size. RF-DETR run records additionally flag uncontrolled online augmentation, preventing an isolated attribution of its differences to offline augmentation.

The missing strong YOLO12-M test entry is retained as a gap. Its validation export and the short probe run in an older generated report are not substitutes. No interpolation across architectures or augmentation strategies is used to complete the table.

\subsection{Effect of MTSD Input Resolution}
\label{sec:ev-resolution}
Table~\ref{tab:ev-resolution} compares the available strong-augmentation test results at 640, 960 and 1280 pixels. Unlike an inference-only resizing experiment, these entries correspond to separately trained configurations at the stated resolution. Effective optimiser batch size is preserved in the controlled higher-resolution YOLO runs through reduced physical batches and gradient accumulation, as described in Subsection~\ref{subsec:mtsd-augmentation-resolution}.

% SOURCE: Results/MTSD-Results/Unified-Evaluation-NoTourist/20260818-084730-pending-yolo-test/*/metrics_summary.json; Results/MTSD-Runs/*/strongaug-*/unified_evaluation/metrics_summary.json.
\begin{table}[htbp]
\centering
\small
\caption{Available strong-augmentation \gls{MTSD} test \gls{mAP}@50--95 (\%) by input resolution. Dashes indicate that a corresponding unified test result was not located, even where a validation result exists.}
\label{tab:ev-resolution}
\renewcommand{\arraystretch}{1.15}
\begin{tabularx}{\textwidth}{@{}>{\raggedright\arraybackslash}Xrrr@{}}
\toprule
Model & 640 & 960 & 1280 \\
\midrule
YOLO11-N & 59.24 & 64.67 & -- \\
YOLO11-S & 62.83 & 68.45 & 72.24 \\
YOLO11-M & 65.36 & 72.15 & 75.34 \\
YOLO12-N & -- & 64.61 & 69.99 \\
YOLO12-S & -- & 67.45 & -- \\
YOLO26-N & 59.66 & 65.53 & -- \\
YOLO26-S & 64.21 & 71.52 & 75.27 \\
YOLO26-M & 67.99 & 73.75 & 76.57 \\
\bottomrule
\end{tabularx}
\end{table}


For YOLO26-M, increasing resolution from 640 to 960 pixels improves \gls{mAP}@50--95 from 67.99\% to 73.75\%, followed by a further improvement to 76.57\% at 1280 pixels. YOLO11-M follows the same pattern, increasing from 65.36\% to 72.15\% and then 75.34\%. These gains exceed the small-to-medium model-scale increments observed on the mild 640-pixel baselines, suggesting that retaining spatial information is at least as important as adding parameters within this experimental setting.

The gains diminish at the highest resolution, while computational cost continues to increase. For YOLO26-M, the stored prediction-loop latency rises from 94.98 to 112.97 and 131.98\,ms/image. The final resolution increment therefore adds 2.82 points of \gls{mAP}@50--95 for approximately 19\,ms/image. YOLO26-S at 1280 pixels reaches 75.27\%, only 1.30 points below YOLO26-M at that resolution, while requiring 115.07\,ms/image in the same evaluator. The operational choice consequently depends on the value assigned to additional small-sign coverage and the available processing budget.

% SOURCE: Results/MTSD-Results/Unified-Evaluation-NoTourist/20260722-all13/*/metrics_summary.json; Results/MTSD-Runs/*/*mtsdqa1aug_img*/unified_evaluation/metrics_summary.json.
\begin{table}[htbp]
\centering
\small
\caption{Available mild-augmentation MTSD test \gls{mAP}@50--95 (\%) by input resolution. The 640-pixel baseline and higher-resolution runs differ in recorded effective batch settings, so gains are descriptive configuration differences.}
\label{tab:ev-resolution-mild}
\renewcommand{\arraystretch}{1.15}
\begin{tabularx}{\textwidth}{@{}>{\raggedright\arraybackslash}Xrrr@{}}
\toprule
Model & 640 & 960 & 1280 \\
\midrule
YOLO11-S & 57.54 & 63.51 & 66.31 \\
YOLO11-M & 60.14 & 65.63 & 68.42 \\
YOLO26-S & 58.98 & -- & 68.15 \\
\bottomrule
\end{tabularx}
\end{table}


The retained mild-augmentation resolution runs show the same broad direction: YOLO11-M increases from 60.14\% at 640 pixels to 65.63\% at 960 and 68.42\% at 1280 pixels, while YOLO26-S increases from 58.98\% to 68.15\% between 640 and 1280 pixels. The earlier 640-pixel runs and these higher-resolution records differ in effective batch settings, so this corroborates the practical value of the higher-resolution configurations without isolating resolution as the only changed factor. For the no-augmentation YOLO26-S experiments at 960 and 1280 pixels, only validation exports were located; they are not inserted into the test comparison.

Higher-resolution RF-DETR-S and -M strong runs record 71.17\% at 640 pixels and 72.66\% at 704 pixels respectively. Their native-resolution strong counterparts lack matching test exports, so these values cannot establish a complete within-strategy RF-DETR resolution curve. The repository's tiled evaluation is outside the principal full-image protocol in Methodology and is not added to the resolution comparison.

% FIGURE PLACEHOLDER: Plot strong-augmentation test mAP@50--95 against 640/960/1280 input resolution for each available YOLO series. Leave gaps for missing test exports. Add a second panel showing matched unified-evaluator latency; do not mix validation or tiled results.

\clearpage
\section{Zero-Shot Prompt-Based Localisation Results}
\label{sec:ev-prompt}
\subsection{MDWD Zero-Shot Localisation}
\label{sec:ev-mdwd-prompt}
Table~\ref{tab:ev-mdwd-prompt-macro} compares the five retained prompt-localisation models over the four class-targeted \gls{MDWD} concepts. The values are taken from the canonical formulations in the completed bounded sensitivity run, which retains the model-specific Cosmos resizing and eight-image worker settings described in Methodology. This avoids using the earlier MDWD Cosmos runtime measurements from a different execution configuration. The corresponding per-prompt \(F_1\) values are reported in Table~\ref{tab:ev-mdwd-prompt-detail}.

% SOURCE: Results/PromptDetect/BatchEvaluation/MDWD/20260821-170445-optimized-v2-bounded-sensitivity/per_prompt_metrics.csv; canonical rows only; arithmetic mean of four exported prompt metrics.
\begin{table}[htbp]
\centering
\small
\caption{Prompt-macro localisation results on the full \gls{MDWD} test split. Predictive measures are percentages; latency is ms per image--prompt pair. Four canonical prompts from the final bounded sensitivity run.}
\label{tab:ev-mdwd-prompt-macro}
\renewcommand{\arraystretch}{1.15}
\begin{tabularx}{\textwidth}{@{}>{\raggedright\arraybackslash}Xrrrrr@{}}
\toprule
Model & $P$ & $R$ & $F_1$ & Matched \gls{IoU} & ms/query \\
\midrule
SAM 3 & 58.43 & 63.80 & 56.76 & 89.64 & 112.3 \\
SAM 3.1 & 58.49 & 64.64 & 56.48 & 89.69 & 111.6 \\
LocateAnything-3B & 23.86 & 77.84 & 33.94 & 88.83 & 1942.6 \\
Cosmos R2-2B & 20.21 & 21.99 & 19.94 & 84.42 & 1478.9 \\
Cosmos R2-8B & 27.87 & 60.54 & 37.07 & 85.08 & 2442.5 \\
\bottomrule
\end{tabularx}
\end{table}


SAM~3 and SAM~3.1 achieve the strongest macro-\(F_1\), at 56.76\% and 56.48\% respectively. Their very small difference provides no clear basis for preferring one checkpoint solely on predictive performance. LocateAnything has the highest macro-recall, at 77.84\%, but its precision of 23.86\% yields a much lower \(F_1\). It therefore retrieves many eligible instances while also producing substantial numbers of non-target predictions. Cosmos Reason2-8B improves considerably over the 2B variant, but remains below SAM in macro-\(F_1\).

% SOURCE: Results/PromptDetect/BatchEvaluation/MDWD/20260821-170445-optimized-v2-bounded-sensitivity/per_prompt_metrics.csv; canonical rows only; arithmetic mean of four exported prompt metrics.
\begin{table}[htbp]
\centering
\small
\caption{Per-prompt $F_1$ (\%) for \gls{MDWD}. LA denotes LocateAnything-3B; C2 and C8 denote Cosmos Reason2-2B and -8B.}
\label{tab:ev-mdwd-prompt-detail}
\renewcommand{\arraystretch}{1.15}
\begin{tabularx}{\textwidth}{@{}>{\raggedright\arraybackslash}Xrrrrr@{}}
\toprule
Prompt & SAM 3 & SAM 3.1 & LA & C2 & C8 \\
\midrule
black garbage bag & 85.16 & 84.12 & 59.22 & 29.82 & 52.22 \\
orange garbage bag & 83.12 & 83.12 & 21.71 & 11.95 & 19.86 \\
white organic-waste bag & 52.69 & 49.82 & 27.31 & 17.93 & 41.42 \\
grey recycling bag & 6.06 & 8.86 & 27.52 & 20.05 & 34.76 \\
\bottomrule
\end{tabularx}
\end{table}


The per-prompt results reveal a strong dependence on the target concept. SAM localises \emph{black garbage bag} and \emph{orange garbage bag} effectively, with \(F_1\) exceeding 83\%, while \emph{grey recycling bag} produces only 6.06\% for SAM~3 and 8.86\% for SAM~3.1. The aggregate weakness therefore does not arise from uniformly poor bounding-box placement. Instead, the correspondence between a generic linguistic concept and the locally annotated collection category is uneven. The recycling query is particularly important because supervised detection of Recyclable Material is comparatively strong, showing that the visual category is learnable even where this wording is ineffective for a pretrained prompt model.

\emph{Other Waste} has no class-targeted entry, consistent with its exclusion from the controlled prompt-family experiment. Broad queries such as \emph{domestic waste} and \emph{garbage} are not included in the headline macro-average. Earlier broad-prompt MDWD exports exist, but a complete bounded-protocol broad block for all five models was not located; their measurements are not inserted into this final comparison.

\subsection{MTSD Zero-Shot Localisation}
\label{sec:ev-mtsd-prompt}
The completed \gls{MTSD} export contains eleven class-targeted and synonym-comparison queries for each model. These queries form the headline macro-average in Table~\ref{tab:ev-mtsd-prompt-macro}; each retains all 747 images, including target-absent scenes. The result is a measure of query-conditioned retrieval, rather than the accuracy of identifying an object already known to be present.

% SOURCE: Results/PromptDetect/BatchEvaluation/MTSD/20260818-134723-optimized-v2/per_prompt_metrics.csv; class-targeted and synonym-comparison macro; broad group separately.
\begin{table}[htbp]
\centering
\small
\caption{Prompt-macro localisation results on the full \gls{MTSD} test split. Predictive measures are percentages; latency is ms per image--prompt pair. Eleven class-targeted and synonym prompts; broad prompts are excluded.}
\label{tab:ev-mtsd-prompt-macro}
\renewcommand{\arraystretch}{1.15}
\begin{tabularx}{\textwidth}{@{}>{\raggedright\arraybackslash}Xrrrrr@{}}
\toprule
Model & $P$ & $R$ & $F_1$ & Matched \gls{IoU} & ms/query \\
\midrule
SAM 3 & 22.71 & 58.04 & 25.74 & 85.51 & 184.3 \\
SAM 3.1 & 22.24 & 58.18 & 25.50 & 85.45 & 189.4 \\
LocateAnything-3B & 5.37 & 56.84 & 9.44 & 94.80 & 6871.3 \\
Cosmos R2-2B & 10.23 & 45.93 & 15.42 & 90.38 & 2602.3 \\
Cosmos R2-8B & 14.63 & 59.07 & 21.78 & 92.01 & 2331.8 \\
\bottomrule
\end{tabularx}
\end{table}


SAM again achieves the highest macro-\(F_1\), although the values of 25.74\% and 25.50\% are substantially below those for waste localisation. The principal limitation is precision: all five models have macro-precision below 23\%. LocateAnything is the clearest example of why matched-box \gls{IoU} alone is insufficient. Its mean matched \gls{IoU} is 94.80\%, but precision is only 5.37\% and \(F_1\) is 9.44\%. The model can place accurate boxes on some eligible signs while failing to restrict its output reliably to the requested semantic category.

% SOURCE: Results/PromptDetect/BatchEvaluation/MTSD/20260818-134723-optimized-v2/per_prompt_metrics.csv; class-targeted and synonym-comparison macro; broad group separately.
\begin{table}[htbp]
\centering
\small
\caption{Per-prompt $F_1$ (\%) for \gls{MTSD}. LA denotes LocateAnything-3B; C2 and C8 denote Cosmos Reason2-2B and -8B.}
\label{tab:ev-mtsd-prompt-detail}
\renewcommand{\arraystretch}{1.15}
\begin{tabularx}{\textwidth}{@{}>{\raggedright\arraybackslash}Xrrrrr@{}}
\toprule
Prompt & SAM 3 & SAM 3.1 & LA & C2 & C8 \\
\midrule
pedestrian crossing sign & 62.22 & 60.74 & 17.67 & 18.16 & 24.32 \\
directional sign & 7.23 & 7.02 & 5.45 & 4.83 & 7.19 \\
back of a traffic sign & 34.89 & 35.40 & 15.50 & 18.75 & 35.25 \\
stop sign & 40.60 & 40.67 & 12.41 & 21.78 & 27.81 \\
no entry sign & 29.20 & 28.50 & 14.19 & 24.89 & 36.96 \\
one way sign & 26.58 & 26.04 & 10.65 & 21.57 & 27.95 \\
roundabout ahead sign & 12.48 & 11.47 & 3.48 & 7.53 & 15.31 \\
no through road sign & 10.94 & 10.65 & 3.75 & 10.64 & 15.63 \\
T-junction sign & 0.00 & 0.00 & 5.89 & 11.33 & 13.39 \\
blind-spot convex mirror & 54.76 & 55.81 & 13.64 & 24.91 & 30.83 \\
street sign & 4.27 & 4.21 & 1.21 & 5.26 & 4.92 \\
\bottomrule
\end{tabularx}
\end{table}


The strongest SAM query is \emph{pedestrian crossing sign}, followed by \emph{blind-spot convex mirror}. Queries for street and directional signs are weak across the model families. The synonym-comparison results also expose semantic limitations: SAM returns \(F_1=0\) for \emph{T-junction sign}, while \emph{no through road sign} produces a non-zero but low score against the mapped category. As established in Methodology, the stored synonym comparisons reflect alternative terminology and should not all be treated as strictly equivalent lexical paraphrases. In particular, \emph{one way sign} is not part of the controlled no-entry paraphrase family.

% SOURCE: Results/PromptDetect/BatchEvaluation/MTSD/20260818-134723-optimized-v2/per_prompt_metrics.csv; class-targeted and synonym-comparison macro; broad group separately.
\begin{table}[htbp]
\centering
\small
\caption{Broad \gls{MTSD} prompts: $F_1$ (\%), reported separately from the headline macro-average.}
\label{tab:ev-mtsd-broad}
\renewcommand{\arraystretch}{1.15}
\begin{tabularx}{\textwidth}{@{}>{\raggedright\arraybackslash}Xrrrrr@{}}
\toprule
Prompt & SAM 3 & SAM 3.1 & LA & C2 & C8 \\
\midrule
traffic sign & 63.82 & 63.42 & 30.07 & 40.23 & 56.88 \\
roadside traffic-control object & 21.31 & 20.70 & 27.96 & 22.91 & 29.59 \\
\bottomrule
\end{tabularx}
\end{table}


Broad retrieval is easier for some models because predictions on several semantic categories become eligible. SAM~3 reaches 63.82\% \(F_1\) for \emph{traffic sign}, compared with its 25.74\% class-targeted and synonym macro-average. This difference supports a distinction between locating sign-like assets and retrieving a specific annotated sign category. It does not represent an improvement on the same task. The much lower result for \emph{roadside traffic-control object} further shows that broader wording is not automatically more effective.

\subsection{Prompt-Sensitivity Analysis}
\label{sec:ev-sensitivity}
The controlled sensitivity experiment retains four formulations for each of four \gls{MDWD} and five \gls{MTSD} target families. Canonical, lexical-paraphrase, descriptive-paraphrase and instruction variants share target mappings and test imagery. Tables~\ref{tab:ev-mdwd-sensitivity} and~\ref{tab:ev-mtsd-sensitivity} summarise within-family variability and prediction consistency rather than comparing unrelated prompts.

% SOURCE: Results/PromptDetect/BatchEvaluation/MDWD/20260821-170445-optimized-v2-bounded-sensitivity/{prompt_sensitivity_per_model,prompt_consistency_per_model}.csv.
\begin{table}[htbp]
\centering
\small
\caption{Controlled MDWD prompt sensitivity. Entries are percentages: mean within-family population standard deviation and range of $F_1$, canonical and instruction $F_1$, matched-\gls{IoU} variability, and prediction-pair agreement.}
\label{tab:ev-mdwd-sensitivity}
\renewcommand{\arraystretch}{1.15}
\begin{tabularx}{\textwidth}{@{}>{\raggedright\arraybackslash}Xrrrrrr@{}}
\toprule
Model & $\sigma_{F_1}$ & $\Delta F_1$ & Canon. & Instr. & $\sigma_{IoU}$ & $F_{1,\mathrm{pair}}$ \\
\midrule
Cosmos R2-2B & 7.35 & 20.37 & 19.94 & 32.58 & 1.59 & 55.15 \\
Cosmos R2-8B & 2.74 & 7.13 & 37.07 & 35.76 & 0.31 & 81.55 \\
LocateAnything-3B & 5.10 & 12.45 & 33.94 & 24.69 & 0.50 & 69.06 \\
SAM 3 & 12.35 & 29.27 & 56.76 & 36.52 & 10.29 & 79.93 \\
SAM 3.1 & 12.04 & 29.42 & 56.48 & 38.11 & 0.69 & 79.12 \\
\bottomrule
\end{tabularx}
\end{table}


On \gls{MDWD}, SAM~3 has a mean within-family \(F_1\) range of 29.27 points and a mean population standard deviation of 12.35 points. Its mean instruction score is 36.52\%, compared with 56.76\% for canonical wording. SAM~3.1 behaves similarly. Cosmos Reason2-2B instead improves from 19.94\% for canonical prompts to 32.58\% for instructions. Thus, a single prompt-writing strategy does not transfer uniformly between a segmentation-based interface and a generative localisation interface. Cosmos Reason2-8B has a smaller mean range, at 7.13 points, although this greater stability accompanies lower localisation accuracy than SAM's strongest formulations.

The stored SAM~3 matched-\gls{IoU} variability on MDWD is unusually large, at 10.29 points, compared with 0.69 for SAM~3.1. This statistic includes the evaluator's zero value when a formulation produces no successful match; it must not be interpreted solely as movement of boxes that remain correctly detected. More generally, both low variability caused by repeated failure and high prediction agreement between empty outputs require interpretation alongside retrieval performance.

% SOURCE: Results/PromptDetect/BatchEvaluation/MTSD/20260822-102455-optimized-v2-bounded-sensitivity/{prompt_sensitivity_per_model,prompt_consistency_per_model}.csv.
\begin{table}[htbp]
\centering
\small
\caption{Controlled MTSD prompt sensitivity. Entries are percentages: mean within-family population standard deviation and range of $F_1$, canonical and instruction $F_1$, matched-\gls{IoU} variability, and prediction-pair agreement.}
\label{tab:ev-mtsd-sensitivity}
\renewcommand{\arraystretch}{1.15}
\begin{tabularx}{\textwidth}{@{}>{\raggedright\arraybackslash}Xrrrrrr@{}}
\toprule
Model & $\sigma_{F_1}$ & $\Delta F_1$ & Canon. & Instr. & $\sigma_{IoU}$ & $F_{1,\mathrm{pair}}$ \\
\midrule
Cosmos R2-2B & 2.96 & 7.42 & 19.45 & 21.60 & 0.86 & 64.39 \\
Cosmos R2-8B & 3.67 & 9.80 & 27.05 & 30.49 & 0.89 & 69.35 \\
LocateAnything-3B & 2.44 & 6.28 & 12.28 & 8.28 & 0.96 & 55.15 \\
SAM 3 & 10.45 & 27.36 & 39.85 & 23.59 & 1.18 & 61.64 \\
SAM 3.1 & 10.49 & 27.32 & 39.44 & 22.74 & 1.23 & 61.17 \\
\bottomrule
\end{tabularx}
\end{table}


On \gls{MTSD}, SAM's canonical formulations perform considerably better than the average over all four formulations. SAM~3 has canonical macro-\(F_1\) of 39.85\% across the five controlled families, but its mean within-family range is 27.36 points. This canonical score differs from the 25.74\% headline result because the controlled study covers five families rather than the eleven-query main inventory. The two quantities should therefore not be interpreted as repeated estimates of the same aggregate.

Prediction consistency provides a complementary assessment. Cosmos Reason2-8B has mean pairwise box \(F_1\) of 81.55\% on MDWD but only 69.35\% on MTSD, while SAM~3 records 79.93\% and 61.64\% respectively. These measures compare output sets without reference to annotation correctness. High agreement may indicate stable localisation, stable semantic confusion or repeated empty predictions. The results therefore favour retaining a documented prompt inventory for reproducibility rather than assuming that fluent reformulation leaves model behaviour unchanged.

% FIGURE PLACEHOLDER: For each controlled target family, show four F1 points (canonical, lexical, descriptive, instruction), grouped by model. Add matched-IoU and pairwise-consistency panels, marking zero-match cases explicitly. Source: final bounded prompt_sensitivity and prompt_consistency CSVs.

\subsection{Supervised Detection versus Zero-Shot Localisation}
\label{sec:ev-paradigms}
The strongest evidence for a paradigm difference is the reliability of semantic targeting. On \gls{MDWD}, supervised YOLO26-L achieves class \(F_1\) of 92.40\% for Mixed Waste and 93.67\% for Recyclable Material in the confidence-0.25 audit. The corresponding SAM~3 canonical queries achieve 85.16\% and 6.06\% at the prompt operating point of 0.30. The second contrast is particularly large, although the different thresholds, output processing and reporting conventions prevent interpreting these differences as a fully harmonised effect size. Orange CMD provides a more favourable case for prompting: its SAM~3 \(F_1\) of 83.12\% is close descriptively to the supervised class score of 85.31\%.

For \gls{MTSD}, the low class-targeted prompt precision contrasts with the high precision available along the supervised confidence curves. Both tracks report standard harmonic-mean \(F_1\), but the detector table uses pooled object counts at the best test-set confidence, whereas the prompt table reports query-macro \(F_1\) at a fixed threshold, with synonym queries receiving separate weight. Subtracting these headline scores would still conflate averaging, target inventory and operating-point differences. A final numerical cross-paradigm ranking on a common target inventory and independently fixed detector threshold is not available in the retained comparison tables. The Tourist Sign ignore-region discrepancy must also be resolved before such a comparison can be described as fully conforming to Methodology.

Within these limits, the evidence supports supervised detection for repeated monitoring of the defined taxonomies, where category-specific training yields reliable retrieval and all classes are predicted in one pass. Prompt-based localisation offers the flexibility to express a target without updating model parameters, and broad traffic-sign retrieval is more successful than specific sign-category queries. That flexibility is accompanied by query-dependent errors and repeated inference cost. It is therefore more immediately suited to exploratory retrieval or analyst-assisted inspection than to replacing the trained detector in the principal benchmark.

\clearpage
\section{Traffic-Sign Attribute Classification Results}
\label{sec:ev-attributes}
\subsection{Overall Representation and Transfer-Strategy Comparison}
\label{sec:ev-attribute-overall}
All 18 principal configurations have retained test metrics for the same 1,890-crop manifest. Table~\ref{tab:ev-attributes} compares their mean macro-\(F_1\), bootstrap intervals and parameter audits. The older six-model results in the July consolidated tables use a different dataset scope and, in some cases, different backbone configurations. They are not combined with the current representation study.

% SOURCE: Scripts/MTSD-Scripts/AttributeClassification/outputs/metrics/*/test_metrics.json; 18 non-smoke configurations, n_images=1890. Numerical confusion matrices used directly.
\begin{table}[htbp]
\centering
\small
\caption{All 18 principal attribute configurations on 1,890 test crops. Mean macro-$F_1$ and its stored 95\% percentile interval are percentages; total and trainable parameter counts are in millions.}
\label{tab:ev-attributes}
\renewcommand{\arraystretch}{1.15}
\begin{tabularx}{\textwidth}{@{}>{\raggedright\arraybackslash}Xlrrrr@{}}
\toprule
Representation & Adaptation & Mean $F_1$ & 95\% CI & Total & Trainable \\
\midrule
DINOv3-B & Frozen & 80.68 & [79.49, 81.88] & 85.67 & 0.010 \\
DINOv3-L & Frozen & 81.62 & [80.50, 82.69] & 303.14 & 0.013 \\
V-JEPA 2.1-B & Frozen & 78.88 & [77.63, 80.03] & 86.84 & 0.010 \\
V-JEPA 2.1-L & Frozen & 80.85 & [79.69, 82.01] & 304.69 & 0.013 \\
ConvNeXt-B & Frozen & 78.16 & [77.07, 79.24] & 87.58 & 0.013 \\
ConvNeXt-L & Frozen & 78.73 & [77.52, 79.91] & 196.25 & 0.020 \\
LingBot-Vision-B & Frozen & 80.15 & [79.00, 81.22] & 85.68 & 0.010 \\
LingBot-Vision-L & Frozen & 81.18 & [80.02, 82.28] & 303.17 & 0.013 \\
DINOv3-B & LoRA & 88.35 & [87.25, 89.30] & 85.97 & 0.305 \\
DINOv3-L & LoRA & 89.59 & [88.60, 90.51] & 303.93 & 0.800 \\
V-JEPA 2.1-B & LoRA & 88.44 & [87.36, 89.46] & 87.14 & 0.305 \\
V-JEPA 2.1-L & LoRA & 90.06 & [89.12, 90.95] & 305.48 & 0.800 \\
ConvNeXt-B & Full FT & 86.84 & [85.70, 87.89] & 87.58 & 87.578 \\
ConvNeXt-L & Full FT & 87.70 & [86.66, 88.65] & 196.25 & 196.247 \\
LingBot-Vision-B & LoRA & 88.54 & [87.46, 89.48] & 85.97 & 0.305 \\
LingBot-Vision-L & LoRA & 89.21 & [88.21, 90.16] & 303.95 & 0.800 \\
ConvNeXt-B & LoRA & 87.97 & [86.84, 88.98] & 89.02 & 1.457 \\
ConvNeXt-L & LoRA & 87.48 & [86.38, 88.44] & 198.41 & 2.186 \\
\bottomrule
\end{tabularx}
\end{table}


V-JEPA~2.1-L with \gls{LoRA} achieves the highest mean macro-\(F_1\), at 90.06\%, followed by DINOv3-L with \gls{LoRA} at 89.59\% and LingBot-Vision-L with \gls{LoRA} at 89.21\%. The 0.47-point lead over DINOv3-L is small relative to the widths of their marginal bootstrap intervals. Since a paired final comparison is unavailable, this is an ordering of observed point estimates rather than evidence that one representation is consistently superior under resampling or retraining.

The more substantial result is the improvement from adapting the representation. Frozen configurations range from 78.16\% to 81.62\%, while \gls{LoRA} configurations range from 87.48\% to 90.06\%. DINOv3-L is the strongest frozen representation, but V-JEPA~2.1-L becomes the highest-scoring adapted configuration. The ordering therefore depends on the transfer strategy and cannot be inferred solely from linear-probe performance.

Increasing capacity is beneficial for the adapted transformer families but not universally. DINOv3, V-JEPA~2.1 and LingBot-Vision all improve from Base to Large with \gls{LoRA}. ConvNeXt instead records 87.97\% for Base and 87.48\% for Large under \gls{LoRA}. Full \gls{FT} reaches 86.84\% and 87.70\% respectively. This non-monotonic pattern reinforces the need to evaluate feasible representation--adaptation combinations rather than assume that larger models or more trainable parameters must improve transfer.

\subsection{Attribute-Level Performance}
\label{sec:ev-attribute-heads}
Table~\ref{tab:ev-attribute-heads} separates the four attribute heads for every configuration. Physical condition is consistently the weakest attribute, while mounting type and geometric shape are learned much more accurately after adaptation. The difference is therefore a stable property of this experiment rather than a failure restricted to one backbone family.

% SOURCE: Scripts/MTSD-Scripts/AttributeClassification/outputs/metrics/*/test_metrics.json; 18 non-smoke configurations, n_images=1890. Numerical confusion matrices used directly.
\begin{table}[htbp]
\centering
\small
\caption{Attribute-level macro-$F_1$ (\%) for the principal test configurations. Each head receives equal weight in the mean reported in Table~\ref{tab:ev-attributes}.}
\label{tab:ev-attribute-heads}
\renewcommand{\arraystretch}{1.15}
\begin{tabularx}{\textwidth}{@{}>{\raggedright\arraybackslash}Xlrrrr@{}}
\toprule
Representation & Adaptation & Angle & Mounting & Condition & Shape \\
\midrule
DINOv3-B & Frozen & 86.86 & 89.85 & 57.67 & 88.34 \\
DINOv3-L & Frozen & 85.91 & 91.27 & 58.35 & 90.95 \\
V-JEPA 2.1-B & Frozen & 85.23 & 86.14 & 55.68 & 88.46 \\
V-JEPA 2.1-L & Frozen & 88.12 & 88.87 & 56.01 & 90.40 \\
ConvNeXt-B & Frozen & 83.70 & 86.51 & 53.39 & 89.03 \\
ConvNeXt-L & Frozen & 83.70 & 88.19 & 55.38 & 87.65 \\
LingBot-Vision-B & Frozen & 87.90 & 89.61 & 53.06 & 90.01 \\
LingBot-Vision-L & Frozen & 87.44 & 89.99 & 55.78 & 91.51 \\
DINOv3-B & LoRA & 92.19 & 95.19 & 69.24 & 96.76 \\
DINOv3-L & LoRA & 92.48 & 95.35 & 72.82 & 97.74 \\
V-JEPA 2.1-B & LoRA & 92.69 & 94.29 & 71.26 & 95.51 \\
V-JEPA 2.1-L & LoRA & 92.97 & 95.26 & 74.70 & 97.32 \\
ConvNeXt-B & Full FT & 91.81 & 93.92 & 65.02 & 96.60 \\
ConvNeXt-L & Full FT & 91.73 & 93.11 & 69.55 & 96.40 \\
LingBot-Vision-B & LoRA & 91.65 & 94.15 & 71.01 & 97.35 \\
LingBot-Vision-L & LoRA & 92.97 & 95.69 & 70.55 & 97.65 \\
ConvNeXt-B & LoRA & 91.90 & 93.28 & 70.22 & 96.47 \\
ConvNeXt-L & LoRA & 91.80 & 92.63 & 68.91 & 96.59 \\
\bottomrule
\end{tabularx}
\end{table}


Viewing-angle macro-\(F_1\) reaches approximately 93\% for the strongest adapted Large models, while mounting type reaches 95.69\% for LingBot-Vision-L with \gls{LoRA}. DINOv3-L obtains 97.74\% for geometric shape, slightly above the other principal contenders. V-JEPA~2.1-L's overall lead is instead associated with physical condition, where it achieves 74.70\%, compared with 72.82\% for DINOv3-L and 70.55\% for LingBot-Vision-L. Equal head weighting ensures that this more difficult task contributes meaningfully to the overall comparison.

% SOURCE: Scripts/MTSD-Scripts/AttributeClassification/outputs/metrics/*/test_metrics.json; 18 non-smoke configurations, n_images=1890. Numerical confusion matrices used directly.
\begin{table}[htbp]
\centering
\small
\caption{Attribute accuracy (\%) for selected adapted representations. Class imbalance makes these values complementary to macro-$F_1$.}
\label{tab:ev-attribute-accuracy}
\renewcommand{\arraystretch}{1.15}
\begin{tabularx}{\textwidth}{@{}>{\raggedright\arraybackslash}Xrrrr@{}}
\toprule
Representation & Angle & Mounting & Condition & Shape \\
\midrule
DINOv3-L & 93.28 & 97.78 & 84.34 & 97.83 \\
V-JEPA 2.1-L & 93.70 & 97.67 & 84.87 & 97.67 \\
LingBot-Vision-L & 93.65 & 97.94 & 80.74 & 97.88 \\
ConvNeXt-B & 92.65 & 96.88 & 84.02 & 96.51 \\
\bottomrule
\end{tabularx}
\end{table}


Accuracy would give a more favourable impression of physical-condition recognition. For V-JEPA~2.1-L with \gls{LoRA}, condition accuracy is 84.87\%, but macro-\(F_1\) is 74.70\%. The discrepancy reflects the dominance of Good signs and the lower performance on minority condition states. Conversely, shape accuracy and macro-\(F_1\) are relatively close for the strongest models, indicating more even performance across that head's labels. These findings support the use of mean macro-\(F_1\) as the primary selection and reporting measure.

\subsection{Per-Class Performance and Confusion Analysis}
\label{sec:ev-attribute-confusion}
Table~\ref{tab:ev-attribute-classes} compares the two highest-scoring representations with \gls{LoRA}. Side views are less accurately recognised than front or back views, while Wall-Mounted signs remain more difficult than Pole-Mounted signs. The latter distinction occurs alongside a large support imbalance: the test set contains 261 wall-mounted crops and 1,629 pole-mounted crops. Such differences remain visible under macro-averaging even when overall accuracy is high.

% SOURCE: Scripts/MTSD-Scripts/AttributeClassification/outputs/metrics/*/test_metrics.json; 18 non-smoke configurations, n_images=1890. Numerical confusion matrices used directly.
\begin{table}[htbp]
\centering
\small
\caption{Per-class test $F_1$ (\%) for V-JEPA~2.1-L and DINOv3-L with LoRA. Support is shared across representations.}
\label{tab:ev-attribute-classes}
\renewcommand{\arraystretch}{1.15}
\begin{tabularx}{\textwidth}{@{}>{\raggedright\arraybackslash}X>{\raggedright\arraybackslash}Xrrr@{}}
\toprule
Attribute & Class & Support & V-JEPA-L & DINOv3-L \\
\midrule
Viewing angle & Front & 699 & 93.15 & 92.71 \\
Viewing angle & Back & 738 & 97.55 & 97.29 \\
Viewing angle & Side & 453 & 88.21 & 87.44 \\
Mounting type & Pole-Mounted & 1629 & 98.64 & 98.71 \\
Mounting type & Wall-Mounted & 261 & 91.88 & 91.98 \\
Physical condition & Good & 1388 & 91.62 & 91.48 \\
Physical condition & Weathered & 392 & 67.00 & 67.45 \\
Physical condition & Heavily Damaged & 110 & 65.49 & 59.51 \\
Geometric shape & Circular & 746 & 98.33 & 98.33 \\
Geometric shape & Quadrangle & 615 & 97.57 & 98.37 \\
Geometric shape & Triangular & 281 & 97.30 & 97.86 \\
Geometric shape & Octagonal & 212 & 96.19 & 94.12 \\
Geometric shape & Pentagon & 36 & 97.22 & 100.00 \\
\bottomrule
\end{tabularx}
\end{table}


Physical condition shows the clearest separation between frequent and infrequent states. Both models achieve approximately 91.5\% \(F_1\) for Good signs, around 67\% for Weathered signs, and lower performance for Heavily Damaged signs. V-JEPA~2.1-L improves the latter class from DINOv3-L's 59.51\% to 65.49\%. Its mean condition advantage therefore reflects, in part, better recognition of severe deterioration rather than an improvement confined to the majority class.

% SOURCE: Scripts/MTSD-Scripts/AttributeClassification/outputs/metrics/*/test_metrics.json; 18 non-smoke configurations, n_images=1890. Numerical confusion matrices used directly.
\begin{table}[htbp]
\centering
\small
\caption{Physical-condition confusion counts for V-JEPA 2.1-L with LoRA. Rows denote \gls{GT}; columns denote predictions.}
\label{tab:ev-condition-vjepa}
\renewcommand{\arraystretch}{1.15}
\begin{tabularx}{\textwidth}{@{}>{\raggedright\arraybackslash}Xrrr@{}}
\toprule
GT / prediction & Good & Weathered & Heavily Damaged \\
\midrule
Good & 1263 & 115 & 10 \\
Weathered & 93 & 267 & 32 \\
Heavily Damaged & 13 & 23 & 74 \\
\bottomrule
\end{tabularx}
\end{table}


% SOURCE: Scripts/MTSD-Scripts/AttributeClassification/outputs/metrics/*/test_metrics.json; 18 non-smoke configurations, n_images=1890. Numerical confusion matrices used directly.
\begin{table}[htbp]
\centering
\small
\caption{Physical-condition confusion counts for DINOv3-L with LoRA. Rows denote \gls{GT}; columns denote predictions.}
\label{tab:ev-condition-dino}
\renewcommand{\arraystretch}{1.15}
\begin{tabularx}{\textwidth}{@{}>{\raggedright\arraybackslash}Xrrr@{}}
\toprule
GT / prediction & Good & Weathered & Heavily Damaged \\
\midrule
Good & 1246 & 133 & 9 \\
Weathered & 80 & 287 & 25 \\
Heavily Damaged & 10 & 39 & 61 \\
\bottomrule
\end{tabularx}
\end{table}


The confusion counts show the operational consequence. V-JEPA~2.1-L correctly recognises 74 of the 110 Heavily Damaged crops, assigning 23 to Weathered and 13 to Good. DINOv3-L correctly recognises 61, assigning 39 to Weathered and ten to Good. Neither model is sufficiently reliable for severe-condition predictions to be interpreted without review. Errors also occur in the opposite direction: V-JEPA~2.1-L assigns 125 Good signs to a deteriorated state, predominantly Weathered. A maintenance workflow must therefore accommodate both missed severe deterioration and unnecessary inspection alerts.

The comparatively low physical-condition annotation stability in Table~\ref{tab:mtsd-annotation-agreement}, with \(\kappa=0.74\), provides relevant context. The model errors are consistent with a task whose severity boundaries require greater interpretation, but annotation agreement is not a model-performance ceiling and does not prove that an individual prediction is correct when it disagrees with the reviewed label. The reviewed annotations remain the evaluation reference.

Geometric shape is more reliable, although support still matters. DINOv3-L correctly classifies all 36 Pentagon examples, yielding a point estimate of 100\% for that class. This small test support does not establish error-free recognition in future imagery. Octagonal signs are more frequently confused, with DINOv3-L's stored matrix showing eight assigned to Circular and four to Quadrangle. Similarly, its viewing-angle matrix contains 55 Side crops assigned to Front, compared with 32 Front crops assigned to Side, indicating an asymmetric boundary between these views.

% FIGURE PLACEHOLDER: Draw row-normalised test confusion matrices for the four heads of V-JEPA 2.1-L LoRA and DINOv3-L LoRA from their numerical test_metrics.json matrices. Annotate counts and class support; retain a common colour scale for matching heads. Do not use the archived 503-crop matrices.

\subsection{Transfer-Strategy Trade-offs}
\label{sec:ev-transfer}
The gains from \gls{LoRA} are large relative to the number of updated parameters. DINOv3-B improves from 80.68\% to 88.35\% mean macro-\(F_1\), while trainable parameters increase from 9,997 head parameters to 304,909 head-and-adapter parameters. V-JEPA~2.1-L improves from 80.85\% to 90.06\% with 799,757 trainable parameters, only 0.2618\% of the adapted model. These results demonstrate effective task-specific adaptation without updating the full representation.

ConvNeXt provides the direct comparison with full \gls{FT}. Its Base \gls{LoRA} configuration attains 87.97\%, compared with 86.84\% for full \gls{FT}, while updating 1.46 million rather than 87.58 million parameters. For Large, full \gls{FT} is slightly higher at 87.70\%, compared with 87.48\% for \gls{LoRA}. Neither difference supports a general claim that full \gls{FT} is the preferred strategy. Moreover, fewer trainable parameters do not guarantee shorter training: ConvNeXt-B \gls{LoRA} records 6,738.4 seconds, compared with 3,877.8 seconds for full \gls{FT}, reflecting different stopping epochs and computation.

The transformer full-\gls{FT} combinations were not evaluated, as stated in Methodology. Their absence is a hardware-feasibility restriction, not evidence that full adaptation would perform worse. Similarly, frozen inference still requires the complete backbone; its very small trainable head is not the size of a deployable stand-alone classifier.

\clearpage
\section{Robustness and Failure Analysis}
\label{sec:ev-robustness}
\subsection{Quantitative Robustness Analysis}
\label{sec:ev-robustness-quantitative}
The retained robustness-slice analysis provides detailed \gls{MDWD} measurements for selected YOLO checkpoints. Table~\ref{tab:ev-robustness} reports YOLO26-L at confidence 0.25 across the seven investigated dimensions. Brightness, contrast and sharpness are divided into empirical terciles of the test imagery; the sharpness measure is an image-level blur proxy rather than an experimental blur intervention. Object-size bins in this audit use the resized 640-pixel coordinate system, unlike the native-coordinate MTSD size analysis.

Image-level slices report pooled \(F_1\), while object-level size, position and frequency slices report recall. Object-level false-positive attribution is not uniquely defined, so these recall values are not converted into invented slice \(F_1\) scores. The audit's minimum support is 15, and the reported slices exceed this requirement. Density groups contain one, two to four, or at least five objects per image. Central objects have box centres within the middle half of both image dimensions; the remaining objects form the near-edge group. Common classes are those with test support at or above the median class count; the complementary group is labelled rare.

% SOURCE: Documents/Final-Reports/Robustness-Slices/20260908-103240/robustness_slice_results.csv and robustness_slice_config.json; yolo26l@YOLO26-EUVIP test only.
\begin{table}[htbp]
\centering
\small
\caption{YOLO26-L \gls{MDWD} robustness slices at confidence 0.25. Support counts images for image-level $F_1$ and objects for object-level recall. Measures are percentages.}
\label{tab:ev-robustness}
\renewcommand{\arraystretch}{1.15}
\begin{tabularx}{\textwidth}{@{}>{\raggedright\arraybackslash}X>{\raggedright\arraybackslash}Xrlr@{}}
\toprule
Dimension & Slice & Support & Measure & Value \\
\midrule
Brightness & high & 123 & $F_1$ & 91.54 \\
Brightness & low & 123 & $F_1$ & 90.53 \\
Brightness & medium & 123 & $F_1$ & 90.86 \\
Contrast & high & 123 & $F_1$ & 88.89 \\
Contrast & low & 123 & $F_1$ & 94.72 \\
Contrast & medium & 123 & $F_1$ & 89.99 \\
Sharpness & blurred & 123 & $F_1$ & 91.98 \\
Sharpness & intermediate & 123 & $F_1$ & 93.69 \\
Sharpness & sharp & 123 & $F_1$ & 87.78 \\
Object density & high clutter & 68 & $F_1$ & 87.72 \\
Object density & low clutter & 159 & $F_1$ & 93.63 \\
Object density & single object & 142 & $F_1$ & 94.81 \\
Object size & large & 594 & Recall & 95.45 \\
Object size & medium & 409 & Recall & 87.78 \\
Object size & small & 103 & Recall & 41.75 \\
Object position & central & 821 & Recall & 93.67 \\
Object position & near edge & 285 & Recall & 70.18 \\
Class frequency & common & 850 & Recall & 86.94 \\
Class frequency & rare & 256 & Recall & 89.84 \\
\bottomrule
\end{tabularx}
\end{table}


The strongest degradation concerns small objects. Recall is 41.75\% for small instances, compared with 87.78\% for medium and 95.45\% for large instances. Near-edge recall is likewise lower than central recall, at 70.18\% versus 93.67\%. Image density also matters: high-clutter scenes have \(F_1=87.72\%\), compared with 94.81\% for single-object scenes. Together, these associations identify limited scale, peripheral position and dense grouping as important conditions for further inspection.

Brightness produces comparatively little separation, with \(F_1\) between 90.53\% and 91.54\%. Contrast and sharpness show a less intuitive pattern: the high-contrast and sharpest terciles perform worse than the low-contrast and least-sharp terciles. These are observational slices, whose scene content and object sizes differ, rather than controlled image degradations. The result does not imply that adding blur or reducing contrast would improve detection. Similarly, the rare-class group has higher recall than the common-class group because it combines Organic Waste with Orange CMD; the aggregate depends on class composition and does not establish that rarity is beneficial.

No corresponding final full-split MTSD image-condition slice report was located. The MTSD size-stratified detector results and current attribute per-class results provide narrower robustness evidence, while the July attribute-slice report and ten-image prompt pilot are historical. They are not presented as quantitative robustness results for the current models.

% FIGURE PLACEHOLDER: Show MDWD YOLO26-L recall by object size and position, and image-level F1 by density, brightness, contrast and sharpness. Include support counts, confidence 0.25, and the coordinate system. Keep object recall and image F1 in separate panels.

\subsection{MDWD Supervised-Detection Failure Cases}
\label{sec:ev-mdwd-failures}
The class and slice results indicate where representative failures should be sought: small peripheral bags, dense overlapping groups and heterogeneous Other Waste instances. High matched-box \gls{IoU} among recovered objects does not offset these misses. Other Waste combines lower precision with lower recall, whereas Orange CMD is more strongly limited by recall, suggesting that a single confidence adjustment will not resolve every class's failure pattern equally.

The numerical evidence supports these priorities but does not establish the visual cause of each error. In particular, overlap, illumination-induced colour ambiguity and partial occlusion should be confirmed from the stored image, annotation and prediction together before being assigned as failure mechanisms.

% FIGURE PLACEHOLDER: Select MDWD test cases from the stored YOLO26-L audit, showing GT and predictions at confidence 0.25: a missed small Orange CMD instance, adjacent overlapping bags, an Other Waste false positive, and a correctly handled comparable scene. Include image identifiers and counts; do not infer failure causes without image review.

\subsection{MTSD Supervised-Detection Failure Cases}
\label{sec:ev-mtsd-failures}
MTSD failures remain concentrated numerically in small-object recovery and the weaker semantic classes. Even the highest-scoring available configuration has small-object AR below 50\%, while Auxiliary Sign and Directional Sign remain substantially below the more distinctive categories. Higher resolution mitigates this limitation but does not eliminate it. A detector used for an inventory must therefore retain a mechanism for revisiting missed or uncertain assets rather than treating a negative detection as proof of absence.

The per-class results cannot determine how much of the remaining loss arises from incorrect category assignment rather than insufficient overlap or a missing prediction. A visual failure analysis should distinguish these cases and use the same native-coordinate test annotations as the unified export. Tourist Sign examples require separate treatment because the stored exports do not implement the prescribed ignore-region behaviour.

% FIGURE PLACEHOLDER: Compare the same MTSD test scenes under strong YOLO26-M at 640 and 1280 pixels. Show small signs recovered only at higher resolution, unresolved Auxiliary/Directional Sign misses, and classification-versus-localisation errors. Mark Tourist Sign ignore regions according to Methodology; do not silently reuse category-deletion scores as corrected measurements.

\subsection{Zero-Shot Localisation Failure Cases}
\label{sec:ev-prompt-failures}
Prompt errors include incomplete retrieval, semantic confusion and formulation-dependent output. For the MTSD \emph{pedestrian crossing sign} query, Cosmos Reason2-2B records 73 \glspl{TP}, 588 \glspl{FP} and 70 \glspl{FN}; 394 false positives overlap a non-target annotated object. Thus, a substantial portion of the false-positive burden is associated with retrieving annotated objects outside the requested category, rather than merely inaccurate geometry on eligible targets. This explains why a model can combine high matched-box \gls{IoU} with weak semantic precision.

Duplicate detections are counted explicitly by the stored matcher, but their contribution must be read from the post-processed outputs rather than presumed from raw model behaviour. The cited pedestrian-crossing query records zero duplicates after post-processing, so its poor precision should not be attributed to duplicate boxes. Prompt-dependent failures are directly supported by the controlled experiments, particularly the large SAM ranges and its failure on some recycling and sign formulations. A useful visual comparison must hold the image and target mapping fixed while changing only the formulation.

% FIGURE PLACEHOLDER: Use final bounded outputs to show a target miss, a semantically related non-target FP, and a fixed-image canonical/paraphrase/instruction comparison. Include a duplicate example only if the stored post-processing metrics identify one. Report the prompt verbatim and identify TP/FP/FN independently of box appearance.

\subsection{Attribute-Classification Failure Cases}
\label{sec:ev-attribute-failures}
Condition severity is the most consequential attribute failure mode. The confusion matrices show substantial movement between adjacent states, particularly Good and Weathered, and between Weathered and Heavily Damaged. Oblique views are also quantitatively less reliable, as reflected in the lower Side-class \(F_1\). These findings justify separate review of condition and viewing-angle uncertainty rather than treating the four attribute outputs as equally dependable.

Occlusion and visually similar surface states are plausible explanations for individual errors, but the retained aggregate results do not quantify their effects independently. The crop-generation rule also excludes the smallest annotated objects, so attribute performance should not be extrapolated to those discarded instances. Since classification uses \gls{GT} crops, further errors caused by inaccurate detector crops remain outside this experiment.

% FIGURE PLACEHOLDER: Pair GT crops and predicted attributes for Good/Weathered and Weathered/Heavily Damaged confusions, plus Front/Side errors. Verify occlusion or ambiguous surface cues visually before describing them. Include correctly classified examples of similar appearance and retain source-image partition identifiers.

\subsection{Exploratory Out-of-Domain Generalisation}
\label{sec:ev-out-of-domain}
No annotated foreign-sign test set or protocol-compatible out-of-domain evaluation was located. Consequently, this chapter makes no quantitative claim about performance on non-Maltese signs. Any future illustrative examples from abroad should remain a separately identified qualitative extension, preferably in an appendix, and should not be used to support international generalisation. The principal conclusions remain restricted to the two Maltese datasets and their observed acquisition conditions.

\clearpage
\section{Computational Efficiency and Deployment Suitability}
\label{sec:ev-deployment}
\subsection{Workstation Inference Performance}
\label{sec:ev-workstation}
The workstation measurements should be interpreted within their respective timing boundaries. Table~\ref{tab:ev-mdwd-runtime} reports native MDWD test timings and retained training durations for the local YOLO configurations. Together with the parameter and checkpoint-size columns in Table~\ref{tab:ev-mdwd-benchmark}, these measurements show that the higher-capacity variants impose measurable computational and storage costs.

% SOURCE: Results/MDWD-Results/{YOLO11,YOLO12,YOLO26}/Model-Size-Comparison/*_summary.csv; test_inference_time_ms and total_training_time_seconds.
\begin{table}[htbp]
\centering
\small
\caption{Native \gls{MDWD} test inference time and training duration. FPS is the reciprocal of the recorded ms/image; it is not an end-to-end municipal processing rate.}
\label{tab:ev-mdwd-runtime}
\renewcommand{\arraystretch}{1.15}
\begin{tabularx}{\textwidth}{@{}>{\raggedright\arraybackslash}Xrrr@{}}
\toprule
Model & ms/image & FPS & Train (h) \\
\midrule
YOLO11-N & 1.11 & 903.7 & 2.89 \\
YOLO11-S & 2.14 & 467.8 & 2.95 \\
YOLO11-M & 3.64 & 274.8 & 3.35 \\
YOLO12-N & 54.70 & 18.3 & 3.35 \\
YOLO12-S & 123.76 & 8.1 & 4.21 \\
YOLO12-M & 280.64 & 3.6 & 4.65 \\
YOLO26-N & 1.55 & 643.5 & 2.97 \\
YOLO26-S & 2.60 & 384.0 & 5.18 \\
YOLO26-M & 3.71 & 269.3 & 5.10 \\
YOLO26-L & 4.81 & 208.0 & 7.39 \\
\bottomrule
\end{tabularx}
\end{table}


% SOURCE: Results/MTSD-Results/Unified-Evaluation-NoTourist/20260818-084730-pending-yolo-test/*/metrics_summary.json; Results/MTSD-Runs/*/strongaug-*/unified_evaluation/metrics_summary.json; inference_latency_ms_per_image and inference_throughput_images_per_second.
\begin{table}[htbp]
\centering
\small
\caption{Strong-augmentation \gls{MTSD} test pipeline timings from the unified evaluator. The timed loop includes image loading and export work; these values are not comparable to native MDWD forward-pass timings.}
\label{tab:ev-mtsd-runtime}
\renewcommand{\arraystretch}{1.15}
\begin{tabularx}{\textwidth}{@{}>{\raggedright\arraybackslash}Xrrrr@{}}
\toprule
Model & Input & \gls{mAP}@50--95 & ms/image & images/s \\
\midrule
YOLO11-S & 640 & 62.83 & 85.64 & 11.68 \\
YOLO11-S & 960 & 68.45 & 98.51 & 10.15 \\
YOLO11-S & 1280 & 72.24 & 114.21 & 8.76 \\
YOLO11-M & 640 & 65.36 & 104.35 & 9.58 \\
YOLO11-M & 960 & 72.15 & 113.54 & 8.81 \\
YOLO11-M & 1280 & 75.34 & 142.29 & 7.03 \\
YOLO26-S & 640 & 64.21 & 86.22 & 11.60 \\
YOLO26-S & 960 & 71.52 & 103.52 & 9.66 \\
YOLO26-S & 1280 & 75.27 & 115.07 & 8.69 \\
YOLO26-M & 640 & 67.99 & 94.98 & 10.53 \\
YOLO26-M & 960 & 73.75 & 112.97 & 8.85 \\
YOLO26-M & 1280 & 76.57 & 131.98 & 7.58 \\
RF-DETR-N & 384 & 63.17 & 183.08 & 5.46 \\
RF-DETR-S & 640 & 71.17 & 225.20 & 4.44 \\
RF-DETR-M & 704 & 72.66 & 284.60 & 3.51 \\
\bottomrule
\end{tabularx}
\end{table}


The MTSD unified prediction-loop values in Table~\ref{tab:ev-mtsd-runtime} are substantially higher than the native MDWD timings, but the difference cannot be assigned solely to dataset difficulty or detector speed because image loading and prediction export are included in the former. Within the strong MTSD runs, YOLO26-S at 1280 pixels records 8.69 images/s, while YOLO26-M records 7.58 images/s. RF-DETR-M at 704 pixels records 3.51 images/s. These are workstation pipeline measurements, not evidence of equivalent throughput on an embedded device.

Prompt latency is reported alongside performance in Tables~\ref{tab:ev-mdwd-prompt-macro} and~\ref{tab:ev-mtsd-prompt-macro}. SAM processes an individual MDWD query in approximately 112\,ms, while the generative models require approximately 1.48--2.44 seconds. On MTSD, LocateAnything requires approximately 6.87 seconds per image--prompt pair. An inventory of multiple queries incurs repeated work; reciprocal single-query latency must not be interpreted as the rate of producing all semantic categories for a municipal scene.

% SOURCE: Scripts/MTSD-Scripts/AttributeClassification/outputs/metrics/*/test_metrics.json; 18 non-smoke configurations, n_images=1890. Numerical confusion matrices used directly.
\begin{table}[htbp]
\centering
\small
\caption{Attribute inference and training costs for selected LoRA configurations. Forward time is amortised per crop at the evaluated batch size; throughput includes the stored end-to-end evaluation pipeline.}
\label{tab:ev-attribute-runtime}
\renewcommand{\arraystretch}{1.15}
\begin{tabularx}{\textwidth}{@{}>{\raggedright\arraybackslash}Xrrr@{}}
\toprule
Representation & Forward ms & Crops/s & Train (min) \\
\midrule
DINOv3-B & 0.86 & 215.37 & 76.94 \\
DINOv3-L & 2.48 & 214.37 & 62.46 \\
V-JEPA 2.1-B & 3.14 & 177.13 & 67.64 \\
V-JEPA 2.1-L & 6.45 & 141.27 & 87.15 \\
LingBot-Vision-B & 1.09 & 210.02 & 74.45 \\
LingBot-Vision-L & 1.91 & 201.05 & 44.52 \\
ConvNeXt-B & 1.09 & 227.01 & 112.31 \\
ConvNeXt-L & 2.13 & 215.84 & 73.96 \\
\bottomrule
\end{tabularx}
\end{table}


Attribute inference is comparatively inexpensive per crop on the workstation, although the forward-pass figures are amortised over batches. DINOv3-B with \gls{LoRA} records 0.86\,ms/crop for the forward pass, compared with 6.45\,ms/crop for V-JEPA~2.1-L with \gls{LoRA}. Their end-to-end rates are 215.37 and 141.27 crops/s respectively. The smaller gap in end-to-end throughput shows that preprocessing and other evaluation work contribute materially to the measured pipeline.

The available MTSD detector exports contain some allocator peak-memory fields, but no complete, uniformly measured memory and serialized-size comparison across all three tracks was located. These fields are not treated as total device-memory requirements. Hosted MDWD RF-DETR training costs are also unavailable, so a common training-time ranking across all architectures cannot be established.

\subsection{Accuracy--Efficiency Trade-offs}
\label{sec:ev-pareto}
A useful operational comparison identifies configurations for which an improvement in predictive performance requires accepting greater computation or storage. Among the MDWD YOLO results, YOLO26-S combines 75.27\% \gls{mAP}@50--95 with a 19.38\,MB checkpoint and 2.60\,ms native inference time. YOLO26-M improves the overlap-averaged score RNA score by only 0.70 points while increasing checkpoint size to 42.00\,MB and latency to 3.71\,ms. YOLO26-L reaches 78.20\% with a 50.54\,MB checkpoint. RF-DETR-M raises standard \(F_1\) from YOLO26-L's 91.80\% to 93.56\% and recall from 86.98\% to 90.69\%, while \gls{mAP}@50--95 remains effectively unchanged at 78.21\%; its 127.54\,MB checkpoint is more than twice as large. Equivalent RF-DETR latency and training-duration measurements were not supplied, so the storage--accuracy trade-off is supported but a common runtime frontier is not.

For the displayed strong MTSD prediction-loop measurements, YOLO26-S at 1280 pixels and YOLO26-M at 1280 pixels both lie on the observed accuracy--latency frontier: the former is faster, while the latter has the highest \gls{mAP}@50--95. YOLO26-M at 960 pixels is also relevant at a slightly lower latency and score. RF-DETR-M's advantage in the mild baseline does not persist as an unconditional deployment recommendation once the higher-resolution strong YOLO configurations are considered. These are empirical comparisons of recorded runs; a single timing measurement does not establish stable dominance across runtime environments.

For attributes, the 0.47-point mean macro-\(F_1\) advantage of V-JEPA~2.1-L over DINOv3-L with \gls{LoRA} accompanies a forward time of 6.45 versus 2.48\,ms/crop. DINOv3-B and LingBot-Vision-B provide further compromises, retaining mean macro-\(F_1\) above 88\% with approximately 86 million total parameters and 0.305 million trainable parameters. The highest-scoring representation is therefore not automatically the most suitable choice when many detected assets must be characterised.

% FIGURE PLACEHOLDER: Three separate accuracy--efficiency plots: native MDWD test mAP@50--95 versus native latency; MTSD test mAP@50--95 versus unified prediction-loop latency, labelled by augmentation and input size; attribute mean macro-F1 versus forward latency with total/trainable parameters distinguished. Add a separate prompt F1-versus-ms/query plot. Identify Pareto-efficient points only within compatible timing groups.

\subsection{Suitability for Edge and Municipal Deployment}
\label{sec:ev-edge}
For recurring kerbside-waste monitoring with a fixed taxonomy, a small supervised YOLO model offers a practical candidate: its checkpoint is compact, its native inference time is low, and its predictive performance is close to larger variants. YOLO26-S is particularly competitive when storage and measured latency matter. Where balanced detection quality is prioritised, RF-DETR-M provides the highest standard \(F_1\), recall and \gls{mAP}@50 in the final MDWD comparison, while YOLO26-L provides nearly identical \gls{mAP}@50--95, the highest precision and a much smaller checkpoint. A deployment choice between these models still requires RF-DETR timing under the same measurement boundary.

Traffic-sign monitoring places greater emphasis on retained spatial detail. Strong YOLO26-S at 1280 pixels offers a useful workstation compromise, while YOLO26-M at 1280 pixels provides the highest available unified test score and better small-object recovery than its lower-resolution counterpart. Either would still miss a substantial fraction of very small signs. For opportunistic street-level monitoring, repeated observations and subsequent human inspection are therefore plausible operational complements, although neither repeated-visit fusion nor its benefit was evaluated here.

Prompt-based models provide semantic flexibility but require care in target specification and substantially more work for a multi-query inventory. Their present results support selective use by an analyst, particularly for broad asset retrieval, rather than unattended replacement of the supervised taxonomy. Attribute \gls{LoRA} models are feasible candidates for a second stage on the measured workstation, but their reported accuracy assumes \gls{GT} crops, and the combined pipeline remains unmeasured.

No matched edge-device latency, energy, quantisation or memory benchmark was retained for the principal comparisons. Deployment suitability is consequently an inference from workstation cost and predictive behaviour, not a demonstrated embedded operating capability. Municipal use would also need a validation-selected operating threshold and resolution of the MTSD ignore-region discrepancy before relying on the reported scores as acceptance criteria.

\clearpage
\section{Synthesis Against the Research Objectives}
\label{sec:ev-objectives}
\subsection{Kerbside-Waste Monitoring}
The available evidence supports supervised category-specific detection for routine kerbside-waste monitoring. RF-DETR-M provides the strongest balanced test result, reaching 93.56\% standard \(F_1\), while YOLO26-L is effectively tied on \gls{mAP}@50--95 and provides higher precision with a substantially smaller checkpoint. The smaller YOLO variants5.27\% retain much of this predictive performance at lower measured cost. Prompt-based localisation is effective for some visually distinctive queries, particularly black and orange bags, but its sensitivity to recycling terminology limits its reliability as a complete alternative.

\subsection{Traffic-Sign Localisation}
Traffic-sign localisation benefits materially from stronger training-data variation and increased input resolution. RF-DETR-M leads the mild baseline, while strong YOLO26-M at 1280 pixels attains the highest available unified \gls{mAP}@50--95. Strong YOLO26-S at the same resolution offers a competitive computational compromise. Small-sign recall and category-specific weakness remain unresolved, and the stored taxonomy treatment qualifies exact compliance with the final evaluation protocol. Prompt models are more successful at broad sign retrieval than precise semantic-category localisation.

\subsection{Traffic-Sign Characterisation}
Representation adaptation is the clearest result in the attribute study. \gls{LoRA} substantially improves all four representation families over frozen probing while updating a small fraction of the backbone parameters. V-JEPA~2.1-L with \gls{LoRA} achieves the highest observed mean macro-\(F_1\), principally through its stronger physical-condition performance, while DINOv3 and LingBot-Vision provide competitive alternatives with lower measured forward cost. Physical condition remains markedly less reliable than geometric shape or mounting type and should be interpreted as a decision-support output rather than an unquestioned maintenance judgement.

\subsection{Overall Municipal-Monitoring Implications}
The findings address the objectives in Section~\ref{sec:aims-objectives} through a conditional rather than a single-metric judgement. Supervised detection is the stronger basis for repeated inventory tasks within a defined taxonomy; prompt localisation offers adaptable queries with less reliable semantic selectivity; and parameter-efficient attribute adaptation supports more detailed characterisation of localised signs. Their respective strengths are complementary, but an integrated end-to-end system was not evaluated. Robustness, observation coverage and human review remain relevant alongside headline accuracy, especially for small objects and condition severity.

\section{Chapter Summary}
\label{sec:ev-summary}
This chapter compared the retained supervised, prompt-localisation and attribute-classification results within the experimental scope established in Methodology. The available supervised measurements show that model scale, augmentation and spatial resolution affect performance differently, with resolution particularly important for small traffic signs. Prompt-based localisation provides useful flexibility but exhibits substantial semantic and formulation sensitivity. Attribute \gls{LoRA} adaptation gives consistent improvements over frozen representations, while physical-condition recognition remains the principal weakness.

The conclusions are bounded by the evidence: the supplied MDWD RF-DETR-S/M results originate from completed Roboflow runs rather than locally retained normalised result files, equivalent RF-DETR timing is unavailable, some MTSD configurations retain only validation results, the stored MTSD taxonomy handling differs from the prescribed ignore-region protocol, and the cross-paradigm \(F_1\) results use different averaging and operating-point conventions. Runtime comparisons preserve their measurement boundaries, and neither edge deployment nor international generalisation is claimed. These qualifications define both the supported findings and the evaluations still required before operational adoption.
